The thesis solves selected combinatorial problems by reducing them to the satisfiability of propositional formulas. The theoretical part introduces the foundations of propositional logic, satisfiability, conjunctive normal form, and constraint satisfaction problems. On this basis, it describes a common reduction procedure in which decision variables and constraints are translated into a Boolean formula and then into a \acrshort{sat} instance. This procedure is applied to selected graph and combinatorial tasks, especially vertex coloring, edge coloring, independent set, and snarks. The implementation part covers Boolean formulas, direct \acrshort{sat} instances, general \acrshort{csp} instances, and graph coloring tasks in detail.

The practical part designs and implements a web application in which the user enters a problem in its domain form, the system creates its \acrshort{sat} encoding, runs an external solver, and interprets the obtained model back into the original task. The design uses domain-driven design, clean architecture, vertical slices, separation of the application service from the background solving service, and asynchronous task processing. The implementation is built on the .NET platform, a React user interface, a PostgreSQL database, Docker containerization, Aspire orchestration, and the CaDiCaL and CryptoMiniSat solvers. The thesis also describes the system interfaces, persistence, local execution, Azure deployment, testing, and measurement limitations. For larger graph inputs, a reliable comparison of the solvers could not be completed because the processing of large text artifacts already failed at the application level.
